\section{Compactness}
\label{sec:fol-metatheory-04}

\begin{theorem}[First-order compactness]
\label{thm:fol-compactness}
If the set of function and relation symbols of \(\Lang\) is not
countable in the sense of \cref{sec:set-prerequisites-05}, assume the
Axiom of Choice in the ambient set theory; if that set is countable,
finite signatures included, no choice principle is assumed.  For every
\(T\subseteq\Sent(\Lang)\) and every \(\sigma\in\Sent(\Lang)\):
\begin{enumerate}
  \item \(T\) has a model iff every finite subset of \(T\) has a model;
        and
  \item if \(T\entails\sigma\), then some finite \(T_0\subseteq T\)
        satisfies \(T_0\entails\sigma\).
\end{enumerate}
Clause 1 is the model form of compactness and clause 2 is the consequence
form; they parallel the two clauses of \cref{thm:prop-compactness}.
\end{theorem}
\begin{proof}
For clause 1, left to right is immediate: a model of \(T\) is a model of
each of its finite subsets.  For right to left, suppose every finite
subset of \(T\) has a model.  If \(T\) were inconsistent, a finite
derivation of \(\bottom\) would use only a finite set
\(T_0\subseteq T\) of assumptions.  Soundness (\cref{thm:fol-soundness})
would make \(T_0\) unsatisfiable, contrary to hypothesis.  Thus \(T\) is
consistent.  If the set of function and relation symbols of \(\Lang\)
is not countable, the Henkin term-structure construction, which uses the
Axiom of Choice through \cref{lem:henkin-lindenbaum}, gives an
\(\Lang^H\)-structure satisfying \(T\), and its \(\Lang\)-reduct is a
model of \(T\) by \cref{lem:reduct-invariance}.  If that set is
countable, a model is obtained without choice.  Were no
\(\Lang\)-structure and assignment to satisfy every member of \(T\), then
\(T\entails\bottom\), and \cref{cor:fol-completeness-countable} would
give \(T\derives_{\ND_{\FOL}}\bottom\), contrary to the consistency of
\(T\).  So some \(\Meta,s\) satisfies every member of \(T\).  The members
of \(T\) are sentences, so by \cref{lem:fol-coincidence} every assignment
in \(\Meta\) satisfies them, and \(\Meta\) is a model of \(T\).

For clause 2, suppose \(T\entails\sigma\).  Completeness
(\cref{thm:fol-completeness}) gives a finite derivation of \(\sigma\)
from \(T\).  Let \(T_0\) be the finite set of assumptions used; soundness
gives \(T_0\entails\sigma\).
\end{proof}

\begin{corollary}
No choice principle is assumed.
The class of finite structures in the empty signature is not axiomatizable by
a first-order theory.
\end{corollary}
\begin{proof}
Suppose \(T\) had exactly the finite nonempty structures as models.  Add the
sentences \(\iota_n\) from \cref{ch:fol-theories}.  Every finite subset of
\(T\cup T_\infty\) has a sufficiently large finite model of \(T\).  The empty signature
is countable, since the empty set injects into \(\N\), so clause 1 of
\cref{thm:fol-compactness} applies in its countable case, which uses no
choice principle, and the union has a model.  Its satisfaction of every \(\iota_n\)
makes it infinite, while its satisfaction of \(T\) makes it finite, a
contradiction.
\end{proof}

